\subsubsection{Linear Structure of Embedding Spaces}\label{sec:linear-structure-of-embedding-spaces}

One of the most remarkable discoveries of Word2Vec is that the embedding space exhibits a \textbf{linear structure}. This means that semantic relationships are not only represented by vector directions, but they can also be \textbf{combined through simple vector arithmetic}.

\highspace
Suppose every word is represented by an embedding vector:
\begin{equation*}
    \mathbf{w} \in \mathbb{R}^m
\end{equation*}
where $m$ is the dimensionality of the embedding space. Since these vectors live in a continuous vector space, we can perform the usual vector operations: addition, subtraction, scalar multiplication. Initially, this might seem like a purely mathematical property. However, after training, these operations acquire a surprising \textbf{linguistic meaning}.

\highspace
\begin{flushleft}
    \textcolor{Green3}{\faIcon{square-root-alt} \textbf{Vector Arithmetic}}
\end{flushleft}
In the previous section, we observed that:
\begin{equation*}
    \mathbf{w}_{\text{Paris}} - \mathbf{w}_{\text{France}} \approx \mathbf{w}_{\text{Rome}} - \mathbf{w}_{\text{Italy}}
\end{equation*}
This tells us that the vector:
\begin{equation*}
    \mathbf{w}_{\text{Paris}} - \mathbf{w}_{\text{France}}
\end{equation*}
Encodes the semantic relation ``\emph{capital of}''. Now suppose we want to answer the following question: \emph{\textbf{what is the capital of Italy?}}. Starting from the known relationship:
\begin{equation*}
    \mathbf{w}_{\text{Paris}} - \mathbf{w}_{\text{France}} \approx \mathbf{w}_{\text{Rome}} - \mathbf{w}_{\text{Italy}}
\end{equation*}
We simply \textbf{apply the same semantic transformation} to the vector of \emph{Italy}:
\begin{equation*}
    \mathbf{w}_{\text{Paris}} - \mathbf{w}_{\text{France}} + \mathbf{w}_{\text{Italy}} \approx \mathbf{w}_{\text{Rome}}
\end{equation*}
This operation can be interpreted as:
\begin{itemize}
    \item Remove the concept of ``\emph{France}'' from the vector of ``\emph{Paris}'';
    \item Keep the concept of ``\emph{capital}'' (the remaining part of the vector);
    \item Add the concept of ``\emph{Italy}'' to the vector.
\end{itemize}
The nearest embedding to the resulting vector is ``\emph{Rome}''.

\highspace
In general, we can express this operation as:
\begin{equation}
    \mathbf{w}_{B} - \mathbf{w}_{A} + \mathbf{w}_{C} \approx \mathbf{w}_{D}
\end{equation}
Where $D$ is the word whose embedding is closest to the resulting vector. Rather than searching for an exact equality, the model searches for the \textbf{nearest neighbor} according to a similarity measure (typically cosine similarity).

\highspace
\begin{flushleft}
    \textcolor{Green3}{\faIcon{question-circle} \textbf{Why does this work?}}
\end{flushleft}
It is important to understand that Word2Vec \textbf{does not contain explicit symbolic knowledge}. It has never been programmed with rules such as:
\begin{itemize}
    \item Paris is the capital of France;
    \item A king is a male monarch;
    \item Copper has chemical symbol Cu.
\end{itemize}
Instead, these relationships emerge because words that participate in similar semantic patterns appear in similar contexts throughout the training corpus. Consequently, \hl{gradient descent gradually arranges the embedding vectors so that many semantic transformations become approximately linear.}

\highspace
\begin{flushleft}
    \textcolor{Red2}{\faIcon{exclamation-triangle} \textbf{Limitations}}
\end{flushleft}
Although these analogies are impressive, they are only \textbf{approximate}. The equations are \textbf{not exact identities}:
\begin{equation*}
    \mathbf{w}_{\text{Paris}} - \mathbf{w}_{\text{France}} + \mathbf{w}_{\text{Italy}} \neq \mathbf{w}_{\text{Rome}}
\end{equation*}
Instead, the resulting vector is simply \textbf{very close} to the embedding of \emph{Rome}. The final answer is obtained by finding the word whose embedding has the highest similarity (typically measured with cosine similarity). Moreover, not every semantic relationship is linear. Some concepts are much more complex and cannot be captured by a single constant vector offset.